\section{Tip-based learning}\label{sec_v4v}

We now consider our second hypothesis: that users learn from and respond to tip feedback by shifting their posting behavior towards better rewarded post types.

\subsection{Model of tip-based learning}

\subsubsection*{Setup}

We first extend the model in Section \ref{sec_pay_to_post_model} to account for user learning over post types. Let there be $K$ post types. Ideas of type $k$ arrive for user $i$ at rate $\gamma_{ik}$, and their quality and intrinsic motivation are distributed according to $f^{ik}(q, \epsilon) = f^{ik}_{\epsilon|q} (\epsilon | q) f^{ik}_q(q)$. As in Section \ref{sec_pay_to_post_model}, we assume that within each post type, $f^{ik}(q,\epsilon)$ satisfies log concavity and the monotone likelihood ratio property.

We parameterize the expected return to a type-$k$ post with:
\begin{align}
V_k(q; \Theta_k) = \Theta_k v_k(q)  \label{eq_type_return}
\end{align}
where $v_k(q)$ is known, continuous, and strictly increasing in $q$. $\Theta_k$ is an unknown, type-specific parameter which reflects how generously posts of type $k$ are rewarded, holding post quality fixed. 

Let $\mathcal{I}_{it}$ denote the user $i$'s information at time $t$, including the tip rewards previously received by the user's posts, and define
\begin{align}
A_{ikt} = E \left[ \Theta_k \mid \mathcal{I}_{it} \right] \label{eq_Mkt}
\end{align}
Because the user is risk-neutral, the perceived value of a type-$k$ idea of quality $q$ is
\begin{align}
\widehat{V}_{ikt}(q) &= E\left[ V_k(q;\Theta_k) \mid \mathcal{I}_{it} \right] \nonumber \\
&= A_{ikt} v_k(q)
\end{align}
The user therefore posts the idea if and only if $A_{ikt} v_k(q) e^{\epsilon} > c_t$, or equivalently:
\begin{align}
\alpha_{ikt} + \ln v_{k}(q) + \epsilon > \kappa_t
\end{align}
where $c_t$ is the posting cost at time $t$, $\kappa_t = \ln c_t$, and $\alpha_{ikt} = \ln A_{ikt}$.

\subsubsection*{Priors and signals}

Let $\theta_k \equiv \ln \Theta_k$, and let $\theta_k \sim N(\mu_0, s_0^2)$ denote the user's prior over $\theta_k$. For every user $i$, let $ikm$ index the $m$th post they made of type $k$, and let $\tau(ikm)$ be the time that the post was made.

After making post $m$ of type $k$, we assume that the user observes the amount of zaps received by the post after 48 hours, $Z_{ikm}$.\footnote{In reality, users can observe zaps as they arrive. However, we abstract away from the full zaps arrival process in order to keep the model parsimonious.}
$Z_{ikm}$ is a noisy estimate of $V_k(q_{ikm} | \Theta_k)$. The noise arises from the inherent randomness of the zaps arrival process, as described in Section \ref{sec_pay_to_post_model}. We assume that the noise can be parameterized as follows:
\begin{align}
\ln Z_{ikm} = \ln V_{k} (q_{ikm} | \Theta_k) + \eta_{ikm}
\end{align}
or, equivalently:
\begin{align}
z_{ikm} = \theta_k + \ln v_{k} (q_{ikm}) + \eta_{ikm}
\end{align}
where $z_{ikm} = \ln Z_{ikm}$ and $\eta_{ikm} \sim N(0, \sigma_{k}^2)$ is the signal noise.

Now define $r_{ikm}$ as a post's quality-adjusted signal:
\begin{align}
r_{ikm} = z_{ikm} - \ln v_k (q_{ikm}) 
\end{align}
%Let $t$ be the start of a week. 
Let $\mathcal{M}_{ikt}$ and $M_{ikt}$ be the set and number of all type-$k$ posts that user $i$ received a signal for through time $t$: 
\begin{align}
\mathcal{M}_{ikt} = \left\{ m: \tau(ikm) < t - \text{48 hours} \right\}, ~ 
~ M_{ikt} = \left| \mathcal{M}_{ikt} \right|
\end{align}
Define $\bar{r}_{ikt}$ as user $i$'s average quality-adjusted signal for type $k$ posts through time $t$:
\begin{align}
\bar{r}_{ikt} = \frac{1}{M_{ikt}} \sum_{m \in \mathcal{M}_{ikt}} r_{ikm}
\end{align}
The learning process has a Gaussian conjugate form, so the user's posterior beliefs at time $t$ are given by:
\begin{align}
\theta_k | \mathcal{I}_{it} \sim N(\mu_{ikt}, s_{ikt}^2)
\end{align}
where
\begin{align}
\mu_{ikt} &= \frac{\sigma_k^2 \mu_0 + s_0^2 M_{ikt} \bar{r}_{ikt}}{\sigma_k^2 + M_{ikt} s_0^2} \\
s_{ikt}^{2} &= \frac{s_0^2 \sigma_{k}^2}{ \sigma_k^2 + M_{ikt} s_0^2}
\end{align}
In words, the user's posterior mean is a weighted average of their past quality-adjusted signals, and the posterior variance is decreasing in the number of observed signals.

\subsubsection*{Behavioral implications}

Noting that:
\begin{align}
\alpha_{ikt} = \mu_{ikt} + \tfrac{1}{2} s_{ikt}^2
\end{align}
we can now see that a user who receives a post idea of type $k$ between times $t$ and $t+\Delta$ will post the idea if and only if:\footnote{For simplicity, we have also assumed that post costs are constant within times $(t, t+\Delta)$.}
\begin{align}
\mu_{ikt} + \tfrac{1}{2} s_{ikt}^2 + \ln v_k (q) + \epsilon > \kappa_t
\end{align}
or, equivalently:
\begin{align}
\epsilon > \kappa_t - \mu_{ikt} - \tfrac{1}{2} s_{ikt}^2 - \ln v_k (q)
\end{align}
Conditional on receiving a post idea, the probability over $(q,\epsilon)$ that the user posts is therefore:
\begin{align}
P_{ikt} = \int S_{\epsilon|q}^{ik}\left(\kappa_t - \mu_{ikt} - \tfrac{1}{2}s_{ikt}^2 - \ln v_k (q) \mid q \right) f^{ik}_q(q)dq
\end{align}
and the instantaneous rate that user $i$ makes a type-$k$ post at time $t$ is $\gamma_{ik} P_{ikt}$.






\subsection{Model of post type choice} \label{sec_v4v_model}

To test this hypothesis, we must first develop some notation for modeling user choice of post type. Consider a user $i$ at time $t$, deciding to make a post from a finite number of post types, $c = 1,\ldots,C$. In theory, the post type could measure any number of different post attributes. In practice, and for parsimony's sake, we limit ourselves to a small number of types based primarily on the length of the post, whether it is a link-only post, and whether it includes hyperlinks and images.

Let $V_{itc}$ be user $i$'s utility for making a post of type $c$ at time $t$. Let $y_{itc}$ equal 1 if the user actually makes a type $c$ post at time $t$, and 0 otherwise. The user chooses the post type that maximizes $V_{itc}$. Thus:
\begin{align}
y_{itc} = 1\left[ V_{itc} \geq V_{itc^\prime} ~ \forall c^\prime=1, \ldots, C \right]
\end{align}
We model the user's utility as follows:
\begin{align}
V_{itc} = \alpha E \left[\ln z_{itc} | I_{itc} \right] + \delta_{c} +\gamma_c \cdot t + \epsilon_{itc}
\end{align}
where $E\left[ \ln z_{itc} | I_{itc} \right]$ is the user's expected value of the log of the sats they would receive on the post, conditional on their information set $I_{itc}$. $\delta_{c}$ is separate intercept for each post type, meant to reflect differences in baseline popularity of each post type.  $\gamma_c \cdot t$ are post type specific linear time trends, which capture changes to the popularity of different post types over time which are not explainable by observed tip feedback. $\epsilon_{itc}$ is an error term which captures idiosyncratic features such as the user's intrinsic desire to make a type-$c$ post in the moment. We model $\epsilon_{itc}$ as a Gumbel distribution, which turns the choice model into a standard conditional logit choice model over $C$ alternatives. Our hypothesis is that $\alpha > 0$; that is, the user's expected rewards for post type $c$ will shift his decision-making towards a greater likelihood of making a type-$c$ post.

\subsection{Post types}

To keep the model parsimonious, we categorize posts into five post types based on the number of words and the presence of images or hyperlinks in the post body. Table \ref{tbl_post_types_summary} shows the five types, their prevalence in the data, the average amount of sats earned in the first 48 hours, and the average number of replies received in the first 48 hours.

The first type of post is a link-only post with no words in the post body. Link-only posts are posts that link to an external URL with no additional text in the post body. Link-only posts earn the least amount of sats, and are generally seen by users as low effort posts because they contain no original content. However, link-only posts are the most common type of post on Stacker News.

The next four types of post are based on whether the post body contains above or below the median number of words, and whether the post body contains images or hyperlinks. Posts with more words tend to earn more sats on average than posts with less words, and posts with images or hyperlinks in the body earn more sats when combined with more words. Posts with more words may be associated with more effort and more original content, which may thereby be rewarded with more tips. Posts with images may be more visually appealing than posts without images, and thus also earn more sats. Hyperlinks in posts with many words are usually associated with linked references and citations in the post body, another sign of quality that may be accordingly rewarded. 

Table \ref{tbl_post_types_summary} shows that posts with an above median number of words and with images or hyperlinks in the body are the best rewarded, by a significant margin. Posts with an above median number of words, but without images or hyperlinks in the body, are the second most rewarded type. The typical earnings of below median length posts are about the same with or without images and hyperlinks. And finally, link-only posts are the least rewarded type of post by a significant margin.




\subsection{Information environment assumptions}

It is not possible for us, as the researchers, to know the information set of each user in each moment that a choice was made. We therefore proceed by making different assumptions about the information environment and showing that the results are robust to the choice of assumption.

To start with, we calculate $E\left[ \ln z_{itc} | I_{itc} \right]$ as follows:
\begin{align}
E\left[ \ln z_{itc} | I_{itc} \right] = \frac{\sum_{i^\prime} \sum_{t^\prime < t} y_{i^\prime t^\prime c} w_{itc,i^\prime t^\prime c} \ln z_{i^\prime t^\prime c} }{\sum_{i^\prime} \sum_{t^\prime < t} y_{i^\prime t^\prime c} w_{itc, i^\prime t^\prime c}}
\end{align}
In other words, we calculate $E\left[ \ln z_{itc} | I_{itc} \right]$ as the weighted average of the realized ln(zaps) for all type $c$ posts made before time $t$ (by any user). Different assumptions about the information set are introduced through the choice of weights, $w_{itc,i^\prime t^\prime c}$, which explicitly depend on both $itc$ (the user making the decision) and $i^\prime t^\prime c$ (the post from which information is being drawn). 

As a baseline, we consider uniform weights as long as the post was made after the user became active. That is, $w_{itc,i^\prime t^\prime c} = 1$ so long as $t^\prime$ is after user $i$'s first post. This assumes that user $i$ observes and draws information from all posts made from the time he became active to the current time.

As a second possibility, we consider weights that are based on recency. In this weighting scheme, $w_{itc,i^\prime t^\prime c} = e^{-\lambda (t - t^\prime)}$, where $\lambda$ is chosen for a half-life of 4 weeks. This information environment assumes that the user pays more attention to recent posts than to old posts, and variation is driven by the timing of the user's posts.

As a third possibility, we consider weights that are based on the proximity of the post to the user's peak activity hours. The idea of this weighting scheme derives from the fact that Stacker News is an international platform, and that different users are most active during different times of the day. For each user $i$, we measure $s_{i}^*$, the average time in the day (measured in seconds from midnight) that the user tends to post. Then, the weights are calculated as $w_{itc, i^\prime t^\prime c} = \exp\left(-\frac{d(s_i^*, t^\prime)^2}{2\sigma^2} \right)$, where $d(s_i^*, t^\prime)$ measures the circular time-of-day distance between $t^\prime$ and $s_i^*$. This information environment assumes that users are more aware of posts made during their peak activity hours.  $\sigma^2$ is chosen to generate a half-life of 3 hours.

\subsection{Results}

Table \ref{tbl_learning} shows the estimation results for this choice model under the aforementioned information assumptions. Columns 1-2 of Table \ref{tbl_learning} show the estimated coefficients under uniform weights, columns 3-4 show the coefficients under recency-based weights, and columns 5-6 shows the coefficients under proximity-to-peak-hours weights.

In the Table, ``Observed Tips'' is equivalent to $E\left[ \ln z_{itc} | I_{itc} \right]$. In all six specifications shown, the coefficient on Observed Tips is estimated to be positive and statistically significant. This confirms our hypothesis that $\alpha>0$; in other words, we find that users shift their posting behavior towards the post types that they observed to be better rewarded. This finding is robust to the various assumptions we made about the information environment.

In addition, we enrich the model by interacting Observed Tips with a measure of the user's experience level, measured by the natural log of the number of posts they've made. This interaction captures a natural feature of learning models which says that learning slows down as the user gets more experienced. When a user has many posts under their belt, they will already have a good sense of what works well for them and what doesn't, and observed tips on other posts are no longer as influential in their decision-making. Consistent with learning theory, we find a negative and statistically significant interaction between Observed Tips and the user's experience level. The finding is robust to our different choices of weighting functions.

Lastly, in columns 2, 4, and 6, we augment the model with a variable called ``User Excess Tips''.  ``User Excess Tips'' is equal to:
\begin{align}
\text{User Excess Tips} &= \frac{\sum_{t^\prime < t} y_{it^\prime c} w_{itc, it^\prime c} \ln z_{it^\prime c}}{\sum_{t^\prime < t} y_{it^\prime c} w_{itc, it^\prime c} } - E\left[ \ln z_{itc} | I_{itc} \right] 
\end{align}
User Excess Tips measures the degree to which the user's own log earnings on posts of type $c$, up to time $t$, have exceeded their estimate based on all observed posts. We expect the coefficient on User Excess Tips to be positive, because a user's own posts are expected to be more salient to their belief formation than other users' posts. A user's own posts could be more salient either because they pay more attention to their own posts, or because their excess earnings relative to other users conveys information about a post-type specific match quality. In either case, a user who earns more tips on a particular post type relative to other users is expected to shift their posting towards that post type. In Table \ref{tbl_learning}, we find a positive and statistically significant coefficient for User Excesss Tips when we use uniform weights or weights based on proximity to the user's peak activity hours. When we use recency-based weights, we find that the effect disappears. This suggests that users are responsive to their historical excess feedback, but not necessarily to recent excess feedback conditional on the observed tips from recent posts. 

\subsection{Discussion}

In this section, we found evidence confirming our hypothesis that users respond to tip-driven feedback by shifting their posting behavior towards posts that appear to be better rewarded. Consistent with learning theory, we also found that new users are more responsive to prior observed tips than seasoned users, and that users are additionally responsive to excess tip feedback on their own posts. The findings were robust to various assumptions about the information environment.

If users shift their posting behavior towards better-rewarded types, then over time we should expect to see the share of those types increase in the aggregate, while the share of poorly rewarded types is expected to decrease. Figure \ref{fig_v4v_q_trend} shows that this is indeed what we see. Link-only posts, the lowest effort and least rewarded type of post, decreased precipitously in the aggregate from the start of the sample period to the end. Non-link-only posts all increased, with above-median length posts with images or hyperlinks increasing in share the most.  

The findings suggest that tip-driven monetary feedback mechanisms could help to shape the overall quality of user submitted content, even if the monetary value of the tips is very small. 

